\chapter{Jacobi Fields}
\label{chap:dc-jacobi-fields}

\section{The Jacobi equation}

\begin{definition}[Jacobi field]
  \label{def:dc-ch5-2-1}
  \lean{Riemannian.Jacobi.exists_jacobiField, Riemannian.Jacobi.jacobiField_eqOn}
  \uses{lem:dc-ch5-2-1-parallel-frame, lem:dc-ch5-2-1-ode, lem:dc-ch5-2-1-frame-reduction, lem:dc-ch5-2-1-frame-existence, lem:dc-ch5-2-1-real-curvature}
  \dcref{ch5:2.1}
  Let $\gamma:[0,a]\to M$ be a geodesic in $M$. A vector field $J$ along $\gamma$ is
  a \emph{Jacobi field} if it satisfies
  $$
  \frac{D^2 J}{dt^2}+R(\gamma'(t),J(t))\gamma'(t)=0.\tag{1}
  $$
  on $[0,a]$.

  If $p\in M$, $v\in T_pM$ lies in the domain of $\exp_p$, and we identify
  $T_v(T_pM)$ with $T_pM$, let $v(s)$ satisfy $v(0)=v$, $v'(0)=w$, and set
  $f(t,s)=\exp_p(tv(s))$. Then the variational field
  $J(t)=\frac{\partial f}{\partial s}(t,0)$ is Jacobi and
  $$
  (d\exp_p)_v w=\frac{\partial f}{\partial s}(1,0),\qquad
  (d\exp_p)_{tv}(tw)=\frac{\partial f}{\partial s}(t,0).
  $$

  A Jacobi field is uniquely determined by $J(0)$ and $\frac{DJ}{dt}(0)$, and
  every pair of initial values occurs. Consequently, if $n=\dim M$, the space
  of Jacobi fields along $\gamma$ has dimension $2n$.
\end{definition}
\begin{proof}
  \sketch
  For each fixed $s$, the curve $t\mapsto f(t,s)$ is geodesic. Commuting the
  covariant derivatives in the variation gives
  $$
  0=\frac{D}{\partial s}\Big(\frac{D}{\partial t}\frac{\partial f}{\partial t}\Big)
  =\frac{D}{\partial t}\frac{D}{\partial t}\frac{\partial f}{\partial s}
  +R\Big(\frac{\partial f}{\partial t},\frac{\partial f}{\partial s}\Big)
    \frac{\partial f}{\partial t},
  $$
  which is (1) for $J=\partial f/\partial s$.

  In a parallel orthonormal frame $e_1(t),\dots,e_n(t)$, write
  
  $$
  J(t)=\sum_i f_i(t)e_i(t),\quad a_{ij}=\langle R(\gamma'(t),e_i(t))\gamma'(t),e_j(t)\rangle,\quad i,j=1,\dots,n=\dim M,
  $$
  
  one has $\frac{D^2 J}{dt^2}=\sum_i f_i''(t)e_i(t)$ and
  $R(\gamma',J)\gamma'=\sum_{ij}f_i a_{ij}e_j$. Therefore equation (1) is equivalent
  to the system
  
  $$
  f_j''(t)+\sum_i a_{ij}(t)f_i(t)=0,\quad j=1,\dots,n,
  $$
  
  Standard existence and uniqueness for this second-order linear system gives
  the claim.
\end{proof}

\begin{lemma}[parallel orthonormal frame along a geodesic, coordinate form]
  \label{lem:dc-ch5-2-1-parallel-frame}
  \lean{Riemannian.Jacobi.exists_parallelOrthoFrame_self, Riemannian.Jacobi.exists_parallelOrthoFrame, Riemannian.Jacobi.exists_chartOrthonormalBasis_self, Riemannian.Jacobi.chartMetricInner_const_of_parallelSol}
  \uses{def:dc-ch2-2-5, lem:dc-ch2-3-1-coord}
  \dcref{ch5:2.1}
  Along a geodesic $\gamma:[0,a]\to M$, every orthonormal basis of
  $T_{\gamma(0)}M$ extends uniquely to a parallel orthonormal frame
  $e_1(t),\dots,e_n(t)$.
\end{lemma}
\begin{proof}

  \sketch
  Extend the initial basis by parallel transport. The inner products
  $\langle e_i,e_j\rangle$ are constant by
  \cref{lem:dc-ch2-3-1-coord}, hence remain equal to $\delta_{ij}$.
\end{proof}

\begin{lemma}[Jacobi system in a parallel frame: existence and uniqueness]
  \label{lem:dc-ch5-2-1-ode}
  \lean{Riemannian.Jacobi.exists_isJacobiPairOn, Riemannian.Jacobi.IsJacobiPairOn.eqOn, Riemannian.Jacobi.IsJacobiPairOn.add, Riemannian.Jacobi.IsJacobiPairOn.smul}
  \leanok
  \dcref{ch5:2.1}
  Let $A:[0,a]\to\operatorname{End}(\mathbb R^n)$ be continuous. For every
  $(f_0,f_1)\in\mathbb R^n\times\mathbb R^n$, the system
  $f''+A(t)f=0$ has a unique solution with $f(0)=f_0$ and $f'(0)=f_1$.
  Its solutions form a $2n$-dimensional vector space.
\end{lemma}

\begin{lemma}[Jacobi reduction in a parallel orthonormal frame]
  \label{lem:dc-ch5-2-1-frame-reduction}
  \lean{Riemannian.Jacobi.covariantDerivCoord_frameCombination, Riemannian.Jacobi.covariantDerivCoord2_frameCombination, Riemannian.Jacobi.chartMetricInner_covariantDerivCoord2_frame, Riemannian.Jacobi.chartMetricInner_map_frameCombination_left, Riemannian.Jacobi.frameJacobiComponent, Riemannian.Jacobi.frameExpansion, Riemannian.Jacobi.covariantDerivCoord2_add_map_frameCombination_expand}
  \uses{lem:dc-ch5-2-1-parallel-frame, prop:dc-ch2-2-2}
  \leanok
  \dcref{ch5:2.1}
  In a parallel orthonormal frame $e_1(t),\dots,e_n(t)$, a field
  $J=\sum_i f_i e_i$ is Jacobi if and only if
  $f_j''+\sum_i a_{ij}f_i=0$ for every $j$, where
  $a_{ij}=\langle R(\gamma',e_i)\gamma',e_j\rangle$.
\end{lemma}
\begin{proof}

  \sketch
  Since $\frac{De_i}{dt}=0$, the Leibniz rule of
  \cref{prop:dc-ch2-2-2} gives
  $$
  \frac{D}{dt}\sum_i f_i e_i=\sum_i f_i'\,e_i,\qquad
  \frac{D^2}{dt^2}\sum_i f_i e_i=\sum_i f_i''\,e_i.
  $$
  Pairing the Jacobi equation with $e_j$ extracts
  $f_j''+\sum_i a_{ij}f_i$. Conversely, expansion in the frame gives
  $\frac{D^2J}{dt^2}+R(\gamma',J)\gamma'=\sum_j\big(f_j''+\sum_i a_{ij}f_i\big)e_j$.
\end{proof}

\begin{lemma}[existence and uniqueness of a Jacobi field by its initial data, frame form]
  \label{lem:dc-ch5-2-1-frame-existence}
  \lean{Riemannian.Jacobi.jacobiCoefOp, Riemannian.Jacobi.exists_jacobiField_frame, Riemannian.Jacobi.jacobiField_frame_eqOn}
  \uses{lem:dc-ch5-2-1-parallel-frame, lem:dc-ch5-2-1-ode, lem:dc-ch5-2-1-frame-reduction}
  \leanok
  \dcref{ch5:2.1}
  For every pair $(J_0,w_0)\in T_{\gamma(0)}M\times T_{\gamma(0)}M$, there is a
  unique Jacobi field $J$ along $\gamma$ such that
  $J(0)=J_0$ and $\frac{DJ}{dt}(0)=w_0$.
\end{lemma}
\begin{proof}
  \leanok
  \uses{lem:dc-ch5-2-1-parallel-frame, lem:dc-ch5-2-1-ode, lem:dc-ch5-2-1-frame-reduction}
  \sketch
  Package the coefficients $a_{ij}(t)$ as a continuous operator field $A(t)$ acting on the
  coefficient vector $(f_i)$ by $(A(t)f)_j=\sum_i a_{ij}(t)f_i$. By \cref{lem:dc-ch5-2-1-ode}
  the second-order linear system $f''+A(t)f=0$ has a solution with initial data
  $f_i(0)=\langle J_0,e_i(0)\rangle$, $f_i'(0)=\langle w_0,e_i(0)\rangle$, unique given that
  data. Put $J=\sum_i f_i e_i$. Pairing the Jacobi operator with the frame
  (\cref{lem:dc-ch5-2-1-frame-reduction}) expands
  $\frac{D^2J}{dt^2}+R(\gamma',J)\gamma'=\sum_j\big(f_j''+\sum_i a_{ij}f_i\big)e_j=0$, so $J$
  is a Jacobi field; the orthonormal expansion gives
  $J(0)=\sum_i\langle J_0,e_i(0)\rangle e_i(0)=J_0$, and the frame being parallel gives the
  initial velocity $w_0$. Uniqueness follows from uniqueness of the scalar system.
\end{proof}

\begin{lemma}[the Jacobi coefficient is the curvature contraction read in the chart]
  \label{lem:dc-ch5-2-1-real-curvature}
  \lean{Riemannian.Jacobi.chartCurvatureCoef, Riemannian.Jacobi.chartCurvatureCoef_contDiffOn, Riemannian.Jacobi.chartCurvatureOp, Riemannian.Jacobi.continuousOn_chartCurvatureOp, Riemannian.Jacobi.exists_jacobiField, Riemannian.Jacobi.jacobiField_eqOn}
  \uses{lem:dc-ch5-2-1-frame-existence, rem:dc-ch4-2-6}
  \dcref{ch5:2.1}
  In coordinates $u$ along $\gamma$, the curvature contraction
  $w\mapsto R(\gamma',w)\gamma'$ has components
  $$
  \big(R(t)\,w\big)^l=\sum_{i,j,k}R^l{}_{ijk}\big(u(t)\big)\,\dot u^i(t)\,w^j\,\dot u^k(t),
  \qquad
  R^l{}_{ijk}=\partial_j\Gamma^l_{ik}-\partial_i\Gamma^l_{jk}
    +\sum_s\big(\Gamma^s_{ik}\Gamma^l_{js}-\Gamma^s_{jk}\Gamma^l_{is}\big),
  $$
  with the coefficients of \cref{rem:dc-ch4-2-6}.
\end{lemma}
\begin{proof}

  \sketch
  These coefficients are smooth, so the operator field is continuous along $\gamma$. Hence
  \cref{lem:dc-ch5-2-1-frame-existence} applies to the intrinsic Jacobi equation.
\end{proof}

\begin{remark}[standard tangential Jacobi fields]
  \label{rem:dc-ch5-2-2}
  \lean{Riemannian.Jacobi.isJacobiFieldAlongOn_velocity, Riemannian.Jacobi.isJacobiFieldAlongOn_smul_velocity}
  \uses{def:dc-ch5-2-1}
  \leanok
  \dcref{ch5:2.2}
  Both $\gamma'$ and $t\gamma'$ are Jacobi fields along $\gamma$. The former is
  parallel and nowhere zero for a nonconstant geodesic, while the latter
  vanishes exactly at $t=0$. Jacobi fields orthogonal to $\gamma'$ will be called
  \emph{normal Jacobi fields}.
\end{remark}
\begin{proof}

  \sketch
  The assertions follow from $\nabla_{\gamma'}\gamma'=0$ and
  $R(\gamma',\gamma')\gamma'=0$.
\end{proof}

\begin{example}[Jacobi fields on manifolds of constant curvature]
  \label{ex:dc-ch5-2-3}
  \lean{Riemannian.Jacobi.isJacobiFieldOn_of_constantCurvature, Riemannian.Jacobi.isJacobiFieldOn_constCurvatureSol_pos, Riemannian.Jacobi.isJacobiFieldOn_constCurvatureSol_zero, Riemannian.Jacobi.isJacobiFieldOn_constCurvatureSol_neg}
  \uses{lem:dc-ch4-3-4, lem:dc-ch5-2-3-const-form, lem:dc-ch5-2-3-jacobi-op}
  \leanok
  \dcref{ch5:2.3}
  Let $M$ have constant sectional curvature $K$, let
  $\gamma:[0,\ell]\to M$ be unit speed, and let $J$ be a normal Jacobi field.
  Then
  $$
  R(\gamma',J)\gamma'=KJ.
  $$
  and the Jacobi equation becomes

  $$
  \frac{D^2 J}{dt^2}+KJ=0.\qquad(2)
  $$

  If $w$ is a parallel unit field orthogonal to $\gamma'$, the solution of (2)
  with $J(0)=0$ and $J'(0)=w(0)$ is
  $$
  J(t)=\begin{cases}\dfrac{\sin(t\sqrt{K})}{\sqrt{K}}\,w(t),&\text{if }K>0,\\[1ex] t\,w(t),&\text{if }K=0,\\[1ex]\dfrac{\sinh(t\sqrt{-K})}{\sqrt{-K}}\,w(t),&\text{if }K<0,\end{cases}
  $$
\end{example}
\begin{proof}
  

  \sketch
  By \cref{lem:dc-ch4-3-4}, for every field $T$ along $\gamma$,
  $$
  \langle R(\gamma',J)\gamma',T\rangle
  =K\big(\langle\gamma',\gamma'\rangle\langle J,T\rangle
  -\langle\gamma',T\rangle\langle J,\gamma'\rangle\big)
  =K\langle J,T\rangle.
  $$
  Thus $R(\gamma',J)\gamma'=KJ$. Since $w$ is parallel, substituting
  $J=h w$ reduces (2) to $h''+Kh=0$, whose normalized solutions are the three
  displayed functions.
\end{proof}

\begin{lemma}[constant curvature: the pointwise $(0,4)$ form, coordinate form]
  \label{lem:dc-ch5-2-3-const-form}
  \lean{Riemannian.Jacobi.curvatureFormAt_isConstantCurvature}
  \uses{lem:dc-ch4-3-4}
  \leanok
  \dcref{ch5:2.3}
  At each $q\in M$, the constant-curvature identity
  $\langle R(X,Y)Z,W\rangle=K_0(\langle X,Z\rangle\langle Y,W\rangle-\langle Y,Z\rangle\langle X,W\rangle)$
  from \cref{lem:dc-ch4-3-4} gives the pointwise $(0,4)$-form
  $\langle R(x,y)z,t\rangle_g=K_0(\langle x,z\rangle\langle y,t\rangle-\langle y,z\rangle\langle x,t\rangle)$
  for all $x,y,z,t\in T_qM$.
\end{lemma}

\begin{lemma}[constant curvature: $R(\gamma',J)\gamma'=K J$, coordinate form]
  \label{lem:dc-ch5-2-3-jacobi-op}
  \lean{Riemannian.Jacobi.chartCurvatureOp_isConstantCurvature}
  \uses{lem:dc-ch5-2-3-const-form, cor:dc-ch5-2-9}
  \leanok
  \dcref{ch5:2.3}
  On a manifold of constant sectional curvature $K_0$, if $J$ is normal to a
  unit-speed geodesic $\gamma$, then
  $R(\gamma',J)\gamma'=K_0J$. Consequently the Jacobi equation is
  $D^2J/dt^2+K_0J=0$.
\end{lemma}
\begin{proof}

  \sketch
  Pair $R(\gamma',J)\gamma'$ with an arbitrary vector and apply
  \cref{lem:dc-ch5-2-3-const-form}; the unit-speed and orthogonality hypotheses
  leave $K_0\langle J,\cdot\rangle$.
\end{proof}

\begin{proposition}[Jacobi fields from exponential-map variations]
  \label{prop:dc-ch5-2-4}
  \lean{Riemannian.Jacobi.expDifferential_eq_jacobiField, Riemannian.Jacobi.IsJacobiFieldAlongOn.eqOn_of_initial}
  \uses{def:dc-ch5-2-1}
  \leanok
  \dcref{ch5:2.4}
  Let $\gamma:[0,a]\to M$ be a geodesic and let $J$ be a Jacobi field along
  $\gamma$ with $J(0)=0$. Put $\frac{DJ}{dt}(0)=w$ and $\gamma'(0)=v$. Using the
  natural identification $T_{av}(T_{\gamma(0)}M)\simeq T_{\gamma(0)}M$, construct a
  curve $v(s)$ in $T_{\gamma(0)}M$ with $v(0)=av$, $v'(0)=aw$. Put
  $f(t,s)=\exp_p(\frac{t}{a}v(s))$,
  $p=\gamma(0)$, and define a Jacobi field $\bar J$ by
  $\bar J(t)=\frac{\partial f}{\partial s}(t,0)$. Then $\bar J=J$ on $[0,a]$.
\end{proposition}
\begin{proof}
  \leanok
  \uses{def:dc-ch5-2-1}
  \sketch
  For $s=0$,
  
  $$
  \frac{D}{dt}\frac{\partial f}{\partial s}=\frac{D}{\partial t}\big((d\exp_p)_{tv}(tw)\big)
  =\frac{D}{\partial t}\big(t(d\exp_p)_{tv}(w)\big)
  =(d\exp_p)_{tv}(w)+t\frac{D}{\partial t}\big((d\exp_p)_{tv}(w)\big).
  $$
  
  Therefore, for $t=0$,
  
  $$
  \frac{D\bar J}{dt}(0)=\frac{D}{\partial t}\frac{\partial f}{\partial s}(0,0)=(d\exp_p)_o(w)=w.
  $$
  
  Since $J(0)=\bar J(0)=0$ and $\frac{DJ}{dt}(0)=\frac{D\bar J}{dt}(0)=w$, we
  conclude from uniqueness that $J=\bar J$.
\end{proof}

\begin{corollary}[exponential differential as a Jacobi field]
  \label{cor:dc-ch5-2-5}
  \lean{Riemannian.Jacobi.expDifferential_eq_jacobiField}
  \uses{prop:dc-ch5-2-4, def:dc-ch5-2-1}
  \leanok
  \dcref{ch5:2.5}
  Let $\gamma:[0,a]\to M$ be a geodesic and put $p=\gamma(0)$. Then a Jacobi field
  $J$ along $\gamma$ with $J(0)=0$ is given by

  $$
  J(t)=(d\exp_p)_{t\gamma'(0)}(tJ'(0)),\quad t\in[0,a].
  $$

\end{corollary}

\begin{remark}[arbitrary initial values and derivative notation]
  \label{rem:dc-ch5-2-6}
  \uses{prop:dc-ch5-2-4}
  \dcref{ch5:2.6}
  The construction of \cref{prop:dc-ch5-2-4} extends to Jacobi fields with
  arbitrary initial value. We henceforth write $\frac{DJ}{dt}=J'$ and
  $\frac{D^2 J}{dt^2}=J''$, etc.
\end{remark}

\begin{lemma}[Taylor expansion of $|J|^2$ in a parallel frame (analytic core)]
  \label{lem:dc-ch5-2-7-taylor-ode}
  \lean{Riemannian.Jacobi.norm_sq_jacobi_isLittleO, Riemannian.Jacobi.norm_sq_jacobi_isLittleO_local}
  \uses{def:dc-ch5-2-1}
  \leanok
  \dcref{ch5:2.7}
  Let $J=\sum_i f_i e_i$ be a Jacobi field in a parallel orthonormal frame,
  and put
  $A(t)=\big(\langle R(\gamma',e_i)\gamma',e_j\rangle\big)_{ij}$.
  If $f(0)=0$ and $w=f'(0)$, then
  $$
  |J(t)|^2=\langle w,w\rangle t^2
  -\tfrac13\langle w,A(0)w\rangle t^4+o(t^4),\qquad t\to0.
  $$
\end{lemma}
\begin{proof}

  \sketch
  The frame equation is $f''=-A(t)f$. For
  $g(t)=\langle f(t),f(t)\rangle$, differentiation at $0$ gives
  $$
  g(0)=g'(0)=g'''(0)=0,\quad
  g''(0)=2\langle w,w\rangle,\quad
  g''''(0)=-8\langle w,A(0)w\rangle.
  $$
  Taylor's theorem with Peano remainder yields the expansion. Moreover,
  $\langle w,A(0)w\rangle
  =\langle R(\gamma'(0),w)\gamma'(0),w\rangle$.
\end{proof}

\begin{lemma}[$C^\infty$ smoothness of the parallel frame and the Jacobi coefficient]
  \label{lem:dc-ch5-2-7-frame-smooth}
  \lean{Riemannian.Jacobi.exists_contDiffOn_parallelOrthoFrame, Riemannian.Jacobi.contDiffOn_infty_jacobiCoefOp}
  \uses{lem:dc-ch5-2-1-parallel-frame, lem:dc-ch5-2-1-real-curvature}
  \leanok
  \dcref{ch5:2.7}
  Along a smooth geodesic in a coordinate chart, the parallel orthonormal frame
  of \cref{lem:dc-ch5-2-1-parallel-frame} may be chosen smooth. Consequently,
  $$
  A(t)=\big(\langle R(\gamma',e_i)\gamma',e_j\rangle\big)_{ij}
  $$
  is smooth in $t$.
\end{lemma}
\begin{proof}

  \sketch
  In coordinates the parallel-transport equation is the linear ODE
  $\dot e_i=-\Gamma(\dot u,e_i)(u)$. Its coefficients are smooth, so its
  solutions are smooth. The assertion for $A$ follows from
  \cref{lem:dc-ch5-2-1-real-curvature}.
\end{proof}

\begin{proposition}[fourth-order norm expansion for exponential variations]
  \label{prop:dc-ch5-2-7}
  \lean{Riemannian.Jacobi.norm_sq_jacobi_frame_isLittleO}
  \uses{lem:dc-ch5-2-7-taylor-ode, lem:dc-ch5-2-7-frame-smooth, def:dc-ch5-2-1}
  \leanok
  \dcref{ch5:2.7}
  Let $p\in M$ and $\gamma:[0,a]\to M$ be a geodesic with $\gamma(0)=p$,
  $\gamma'(0)=v$. Let $w\in T_v(T_p M)$ with $|w|=1$ and let $J$ be a Jacobi field
  along $\gamma$ given by
  $$
  J(t)=(d\exp_p)_{tv}(tw),\quad 0\le t\le a.
  $$
  Then the Taylor expansion of $|J(t)|^2$ about $t=0$ is given by
  $$
  |J(t)|^2=t^2-\tfrac{1}{3}\langle R(v,w)v,w\rangle t^4+R(t),\qquad(3)
  $$
  where $\lim_{t\to 0}\frac{R(t)}{t^4}=0$.
\end{proposition}
\begin{proof}
  \leanok
  \uses{lem:dc-ch5-2-7-taylor-ode, lem:dc-ch5-2-7-frame-smooth}
  \sketch
  Since $J(0)=0$ and $J'(0)=w$, differentiation gives
  
  $$
  \langle J,J\rangle(0)=0,\quad\langle J,J\rangle'(0)=2\langle J,J'\rangle(0)=0,\quad\langle J,J\rangle''(0)=2\langle J',J'\rangle(0)+2\langle J'',J\rangle(0)=2.
  $$
  
  Since $J''(0)=-R(\gamma',J)\gamma'(0)=0$, we have
  $\langle J,J\rangle'''(0)=6\langle J',J''\rangle(0)+2\langle J''',J\rangle(0)=0$. Now
  we need the fact
  
  $$
  \nabla_{\gamma'}(R(\gamma',J)\gamma')(0)=R(\gamma',J')\gamma'(0).\qquad(4)
  $$
  
  For any field $W$, at $t=0$,
  
  $$
  \Big\langle\frac{D}{dt}(R(\gamma',J)\gamma'),W\Big\rangle
  =\frac{d}{dt}\langle R(\gamma',W)\gamma',J\rangle-\langle R(\gamma',J)\gamma',W'\rangle
  =\Big\langle\frac{D}{dt}(R(\gamma',W)\gamma'),J\Big\rangle+\langle R(\gamma',W)\gamma',J'\rangle
  =\langle R(\gamma',J')\gamma',W\rangle,
  $$
  
  which proves (4). The Jacobi equation then gives
  $J'''(0)=-R(\gamma',J')\gamma'(0)$. Therefore,
  
  $$
  \langle J,J\rangle''''(0)=8\langle J',J'''\rangle(0)+6\langle J'',J''\rangle(0)+2\langle J'''',J\rangle(0)=-8\langle J',R(\gamma',J')\gamma'\rangle(0)=-8\langle R(v,w)v,w\rangle.
  $$
  
  Taylor's theorem now gives (3).
\end{proof}

\begin{remark}[curvature derivative identity at the initial point]
  \label{rem:dc-ch5-2-8}
  \dcref{ch5:2.8}
  Identity (4) also follows by covariantly differentiating the curvature
  tensor along $\gamma$ and using $J(0)=0$.
\end{remark}

\begin{corollary}[sectional-curvature form of the norm expansion]
  \label{cor:dc-ch5-2-9}
  \lean{Riemannian.Jacobi.chartMetricInner_chartCurvatureOp_eq_sectionalCurvature, Riemannian.Jacobi.chartMetricInner_chartCurvatureOp_eq_curvatureFormAt}
  \uses{def:dc-ch4-3-2, prop:dc-ch5-2-7}
  \leanok
  \dcref{ch5:2.9}
  If $\gamma:[0,\ell]\to M$ is parametrized by arc length (i.e. $|v|=1$) and
  $\langle w,v\rangle=0$, the expression $\langle R(v,w)v,w\rangle$ is the sectional
  curvature at $p$ with respect to the plane $\sigma$ generated by $v$ and $w$.
  Therefore, in this situation,

  $$
  |J(t)|^2=t^2-\tfrac{1}{3}K(p,\sigma)t^4+R(t).\qquad(5)
  $$

\end{corollary}

\begin{lemma}[square root of the $|J|^2$ expansion (analytic core)]
  \label{lem:dc-ch5-2-10-sqrt}
  \lean{Riemannian.Jacobi.sqrt_isLittleO_of_sq_isLittleO}
  \uses{cor:dc-ch5-2-9}
  \leanok
  \dcref{ch5:2.10}
  Let $g\ge0$ satisfy $g(t)=t^2-ct^4+o(t^4)$ as $t\to0$. Then
  $$
  \sqrt{g(t)}=t-\tfrac{c}{2}\,t^3+o(t^3),\qquad t\to 0^+.
  $$
  The one-sided limit is necessary because $\sqrt{t^2}=|t|$.
\end{lemma}
\begin{proof}

  \sketch
  Put $p(t)=t-\tfrac c2t^3$. Then $g-p^2=o(t^4)$ and
  $\sqrt g-p=(g-p^2)/(\sqrt g+p)$. For sufficiently small $t>0$, the denominator
  is at least $t/2$, so the quotient is $o(t^3)$.
\end{proof}

\begin{corollary}[norm expansion and sectional curvature]
  \label{cor:dc-ch5-2-10}
  \lean{Riemannian.Jacobi.norm_jacobi_frame_isLittleO}
  \uses{prop:dc-ch5-2-7, lem:dc-ch5-2-10-sqrt, cor:dc-ch5-2-9}
  \leanok
  \dcref{ch5:2.10}
  With the same conditions as in the previous corollary,

  $$
  |J(t)|=t-\tfrac{1}{6}K(p,\sigma)t^3+\tilde R(t),\quad\lim_{t\to 0}\frac{\tilde R(t)}{t^3}=0.\qquad(6)
  $$

\end{corollary}

\begin{remark}[curvature controls radial spreading]
  \label{rem:dc-ch5-2-11}
  \dcref{ch5:2.11}
  Let
  $$
  f(t,s)=\exp_p tv(s),\quad t\in[0,\delta],\quad s\in(-\varepsilon,\varepsilon),
  $$
  where $|v(s)|=1$, $v(0)=v$, $v'(0)=w$, and $\delta>0$ is chosen so that
  $\exp_p(tv(s))$ is defined. The variation of the radial lines in $T_pM$ has
  norm
  $$
  \Big|\Big(\frac{\partial}{\partial s}tv(s)\Big)(0)\Big|=|tw|=t.
  $$
  By (6), the norm of the corresponding geodesic variation is
  $t-\frac16K(p,\sigma)t^3+o(t^3)$. Thus positive sectional curvature decreases,
  and negative sectional curvature increases, the infinitesimal radial
  separation relative to $T_pM$.
\end{remark}

\section{Conjugate points}

\begin{definition}[conjugate point, multiplicity]
  \label{def:dc-ch5-3-1}
  \lean{Riemannian.Jacobi.IsConjugatePointAt, Riemannian.Jacobi.IsGeodesicConjugatePointAt}
  \uses{def:dc-ch5-2-1}
  \leanok
  \dcref{ch5:3.1}
  Let $\gamma:[0,a]\to M$ be a geodesic and $t_0\in(0,a]$. The point
  $\gamma(t_0)$ is \emph{conjugate} to $\gamma(0)$ along $\gamma$ if some
  nonzero Jacobi field $J$ satisfies $J(0)=J(t_0)=0$. Its \emph{multiplicity}
  is the dimension of the space of such fields. Conjugacy along a geodesic
  segment is symmetric under reversal of the segment.
\end{definition}

\begin{remark}[dimension bound for conjugate multiplicity]
  \label{rem:dc-ch5-3-2}
  \lean{Riemannian.Jacobi.conjugateMultiplicity, Riemannian.Jacobi.conjugateMultiplicity_le_finrank, Riemannian.Jacobi.conjugateMultiplicity_le_finrank_sub_one}
  \uses{rem:dc-ch5-2-2}
  \dcref{ch5:3.2}
  If $\dim M=n$, the Jacobi fields along $\gamma$ that vanish at $0$ form an
  $n$-dimensional space. A conjugate point along a nonconstant geodesic has
  multiplicity at most $n-1$.
\end{remark}
\begin{proof}

  \sketch
  The initial-derivative map $J\mapsto J'(0)$ identifies this space with
  $T_{\gamma(0)}M$; in particular, fields $J_i$ vanishing at $0$ are linearly
  independent exactly when the vectors $J_i'(0)$ are. At a positive endpoint,
  the field $t\gamma'(t)$ of \cref{rem:dc-ch5-2-2} does not vanish, so the kernel
  of endpoint evaluation is a proper subspace and has dimension at most $n-1$.
\end{proof}

\begin{example}[conjugate points in positive constant curvature]
  \label{ex:dc-ch5-3-3}
  \lean{Riemannian.Jacobi.isConjugatePointAt_constCurvature_pos, Riemannian.Jacobi.conjugateMultiplicity_constCurvature_pos_eq}
  \uses{ex:dc-ch5-2-3, def:dc-ch5-3-1, rem:dc-ch5-3-2, prop:dc-ch5-3-5}
  \leanok
  \dcref{ch5:3.3}
  Let $M^n$ have constant sectional curvature $K_0>0$, and let
  $\gamma:[0,\pi/\sqrt{K_0}]\to M$ be unit speed. Then
  $\gamma(\pi/\sqrt{K_0})$ is conjugate to $\gamma(0)$ with multiplicity $n-1$.
  In particular, on the unit sphere $S^n$, the antipode $\gamma(\pi)$ has
  multiplicity $n-1$ along every unit-speed geodesic from $\gamma(0)$.
\end{example}
\begin{proof}

  \sketch
  For each $w_0\perp\gamma'(0)$, let $w$ be its parallel translate along
  $\gamma$. By \cref{ex:dc-ch5-2-3},
  $$
  J(t)=\frac{\sin(t\sqrt{K_0})}{\sqrt{K_0}}\,w(t)
  $$
  is Jacobi and vanishes at both endpoints. These fields form an
  $(n-1)$-dimensional family, while \cref{rem:dc-ch5-3-2} gives the reverse
  bound on multiplicity.
\end{proof}

\begin{definition}[conjugate locus]
  \label{def:dc-ch5-3-4}
  \lean{Riemannian.Jacobi.IsFirstConjugatePointAt, Riemannian.Jacobi.conjugateLocus}
  \uses{def:dc-ch5-3-1}
  \dcref{ch5:3.4}
  The \emph{conjugate locus} $C(p)$ of $p\in M$ is the set of first conjugate
  points along all geodesics starting at $p$. For the unit sphere,
  $C(p)=\{-p\}$.
\end{definition}

\begin{proposition}[conjugate points and critical points of the exponential map]
  \label{prop:dc-ch5-3-5}
  \lean{Riemannian.Jacobi.expDifferential_injective_iff_not_conjugate, Riemannian.Jacobi.injective_jacobiEndpointOfVel_iff_not_conjugate, Riemannian.Jacobi.conjugateMultiplicity, Riemannian.Jacobi.isConjugatePointAt_iff_conjugateMultiplicity_pos, Riemannian.Jacobi.conjugateMultiplicity_eq_finrank_ker_expDifferential}
  \uses{cor:dc-ch5-2-5, def:dc-ch5-3-1}
  \leanok
  \dcref{ch5:3.5}
  Let $\gamma:[0,a]\to M$ be a geodesic with $\gamma(0)=p$, and let
  $t_0\in(0,a]$. Then $q=\gamma(t_0)$ is conjugate to $p$ along $\gamma$ if and
  only if $v_0=t_0\gamma'(0)$ is a critical point of $\exp_p$. Moreover, the
  multiplicity of $q$ equals $\dim\ker(d\exp_p)_{v_0}$.
\end{proposition}
\begin{proof}
  \sketch
  Put $v=\gamma'(0)$. By \cref{cor:dc-ch5-2-5}, every Jacobi field with
  $J(0)=0$ is determined by $w=J'(0)$ and satisfies
  $$
  J(t_0)=(d\exp_p)_{t_0v}(t_0w).
  $$
  Since $t_0\ne0$, nonzero initial derivatives in the kernel correspond exactly
  to nonzero Jacobi fields vanishing at both endpoints. The map
  $w\mapsto t_0w$ is an isomorphism, so it also identifies their space with
  $\ker(d\exp_p)_{t_0v}$ and preserves dimension.
\end{proof}

\begin{proposition}[affine tangential component of a Jacobi field]
  \label{prop:dc-ch5-3-6}
  \lean{Riemannian.Jacobi.chartMetricInner_jacobi_velocity_affine}
  \uses{def:dc-ch5-2-1}
  \leanok
  \dcref{ch5:3.6}
  Let $J$ be a Jacobi field along the geodesic $\gamma:[0,a]\to M$. Then

  $$
  \langle J(t),\gamma'(t)\rangle=\langle J'(0),\gamma'(0)\rangle t+\langle J(0),\gamma'(0)\rangle,\quad t\in[0,a].
  $$

\end{proposition}
\begin{proof}
  \leanok
  \uses{def:dc-ch5-2-1}
  \sketch
  Metric compatibility, the geodesic equation, the Jacobi equation, and the
  curvature symmetries give
  $$
  \langle J',\gamma'\rangle'=\langle J'',\gamma'\rangle=-\langle R(\gamma',J)\gamma',\gamma'\rangle=0.
  $$
  Hence $\langle J',\gamma'\rangle$ is constant, and
  $$
  \langle J,\gamma'\rangle'=\langle J',\gamma'\rangle=\langle J'(0),\gamma'(0)\rangle.
  $$
  Integration yields the stated affine formula.
\end{proof}

\begin{corollary}[constancy of the tangential component]
  \label{cor:dc-ch5-3-7}
  \lean{Riemannian.Jacobi.chartMetricInner_jacobi_velocity_const_of_eq}
  \uses{prop:dc-ch5-3-6}
  \leanok
  \dcref{ch5:3.7}
  If $\langle J,\gamma'\rangle(t_1)=\langle J,\gamma'\rangle(t_2)$, $t_1,t_2\in[0,a]$,
  $t_1\neq t_2$, then $\langle J,\gamma'\rangle$ does not depend on $t$; in
  particular, if $J(0)=J(a)=0$, then $\langle J,\gamma'\rangle(t)\equiv 0$.
\end{corollary}
\begin{proof}

  \sketch
  An affine function taking the same value at two distinct times has zero
  slope. If $J$ vanishes at both endpoints, that constant value is zero.
\end{proof}

\begin{corollary}[normal Jacobi fields with zero initial value]
  \label{cor:dc-ch5-3-8}
  \lean{Riemannian.Jacobi.chartMetricInner_jacobi_velocity_eq_zero_iff, Riemannian.Jacobi.velocityFunctional, Riemannian.Jacobi.finrank_velocityPerp_eq}
  \uses{prop:dc-ch5-3-6, rem:dc-ch5-3-2}
  \leanok
  \dcref{ch5:3.8}
  Suppose $J(0)=0$. Then
  $\langle J'(0),\gamma'(0)\rangle=0$ if and only if
  $\langle J(t),\gamma'(t)\rangle=0$ for every $t$. If $\dim M=n$ and $\gamma$
  is nonconstant, the space of such Jacobi fields has dimension $n-1$.
\end{corollary}
\begin{proof}

  \sketch
  By \cref{prop:dc-ch5-3-6}, the pairing is
  $t\langle J'(0),\gamma'(0)\rangle$. Under the identification
  $J\mapsto J'(0)$, the normal fields form the hyperplane
  $\gamma'(0)^\perp\subset T_{\gamma(0)}M$.
\end{proof}

\begin{proposition}[Jacobi field with prescribed nonconjugate endpoints]
  \label{prop:dc-ch5-3-9}
  \lean{Riemannian.Jacobi.exists_unique_jacobi_of_endpoints}
  \uses{cor:dc-ch5-2-5, def:dc-ch5-3-1}
  \leanok
  \dcref{ch5:3.9}
  Let $\gamma:[0,a]\to M$ be a geodesic, let
  $V_1\in T_{\gamma(0)}M$, and let $V_2\in T_{\gamma(a)}M$. If $\gamma(a)$ is
  not conjugate to $\gamma(0)$ along $\gamma$, there is a unique Jacobi field
  $J$ such that $J(0)=V_1$ and $J(a)=V_2$.
\end{proposition}
\begin{proof}
  \leanok
  \uses{cor:dc-ch5-2-5, def:dc-ch5-3-1}
  \sketch
  For $w\in T_{\gamma(0)}M$, let $J_w$ be the Jacobi field with
  $J_w(0)=0$ and $J_w'(0)=w$. The endpoint map
  $$
  \Theta:T_{\gamma(0)}M\longrightarrow T_{\gamma(a)}M,\qquad
  \Theta(w)=J_w(a),
  $$
  is injective by nonconjugacy, hence is an isomorphism because both spaces
  have dimension $\dim M$. Let $J_0$ be the Jacobi field with initial data
  $(V_1,0)$, choose $w$ with $\Theta(w)=V_2-J_0(a)$, and set $J=J_0+J_w$.
  The difference of two solutions lies in $\ker\Theta$, proving uniqueness.
\end{proof}

\begin{corollary}[normal Jacobi fields form an endpoint basis]
  \label{cor:dc-ch5-3-10}
  \lean{Riemannian.Jacobi.jacobiConjugateBasis, Riemannian.Jacobi.jacobiConjugateEquiv, Riemannian.Jacobi.jacobiEndpointOfVel_map_velocityPerp_eq, Riemannian.Jacobi.metricInner_jacobiJ_velocity_eq_zero}
  \uses{cor:dc-ch5-3-8, prop:dc-ch5-3-5, prop:dc-ch5-3-6, def:dc-ch5-3-1}
  \leanok
  \dcref{ch5:3.10}
  Let $\gamma:[0,a]\to M$ be a geodesic in $M$, $\dim M=n$, and let
  $\mathcal{J}^\perp$ be the space of Jacobi fields with $J(0)=0$,
  $J'(0)\perp\gamma'(0)$. Let $\{J_1,\dots,J_{n-1}\}$ be a basis of
  $\mathcal{J}^\perp$. If $\gamma(t)$, $t\in(0,a]$, is not conjugate to $\gamma(0)$,
  then $\{J_1(t),\dots,J_{n-1}(t)\}$ is a basis for the orthogonal complement
  $\{\gamma'(t)\}^\perp\subset T_{\gamma(t)}M$ of $\gamma'(t)$.
\end{corollary}
\begin{proof}

  \sketch
  By \cref{cor:dc-ch5-3-8}, evaluation at $t$ sends
  $\mathcal J^\perp$ into $\gamma'(t)^\perp$. Nonconjugacy makes this endpoint
  map injective by \cref{prop:dc-ch5-3-5}. Both spaces have dimension $n-1$,
  so it is an isomorphism and carries any basis of $\mathcal J^\perp$ to the
  displayed basis.
\end{proof}

\begin{example}[conjugate loci of the sphere and ellipsoid]
  \label{ex:dc-ch5-3-4}
  \dcref{ch5:3.4}
  \uses{def:dc-ch5-3-4}
  \notready
  On $S^n$, $C(p)=\{-p\}$, for all $p$. The case of $S^n$, however, is not typical.
  A more typical example is given by the ellipsoid, where $C(p)$ is, in general, a
  curve with four singular points (see Fig. 4 of Chap. 13; cf. Braunmühl, A.,
  "Geodätische Linien auf dreiachsigen Flächen 2-Grades", Math. Ann. 20 (1882),
  557-586).
\end{example}
